\chapter{Fourier Pricing and Calibration Engineering}\label{ch:dv:fourier-pricing-and-calibration-engineering}

Every morning a job fits a \omterm{def:dv:stochastic-volatility:heston}{Heston model} to twenty option quotes. On day 58 of this chapter's synthetic market the fitted \omterm{def:dv:stochastic-volatility:sv}{volatility of
volatility} is 0.51; on day 59 it is 0.28, with the fit error unchanged at 0.31 volatility point. Nothing happened in the market that the
quotes can see: the move is noise, amplified by a direction in parameter space along which the fit barely changes. But the exotics priced
with the model see it. The convexity adjustment of a one-year \omterm{def:dv:variance-swaps-and-volatility-derivatives:volswap}{volatility swap} falls from 1.32 to 0.39 volatility point overnight, and a
desk that marks its book this way books a profit or loss that no hedge can explain. This chapter covers the two halves of the job: pricing
fast enough from a characteristic function to calibrate on it, and calibrating in a way whose output a desk can use from one day to the
next.

\section{Pricing by transform}

Chapters 10 and 13 wrote every model through the characteristic function $\varphi(u)=\E[e^{iux_T}]$ of $x_T=\ln(S_T/F)$, and priced with
one integral. That integral comes from Parseval's identity.

\begin{definition}[Lewis formula]\label{def:dv:fourier-pricing-and-calibration-engineering:lewis}
The \emph{Lewis formula}\index{Lewis formula} prices a European call from the characteristic function of $x_T$ along the line
$\operatorname{Im}u=-\frac12$: with $k=\ln(K/F)$,
\[
C=P(0,T)\Bigl(F-\frac{\sqrt{FK}}{\pi}\int_0^\infty\operatorname{Re}\bigl(e^{-iuk}\varphi(u-\tfrac i2)\bigr)\frac{du}{u^2+\frac14}\Bigr).
\]
\end{definition}

\begin{proposition}[Derivation]\label{prop:dv:fourier-pricing-and-calibration-engineering:lewis}
With $g(x)=\min(e^x,e^k)$, $F-C/P(0,T)=F\,\E[g(x_T)]$. For $0<\operatorname{Im}z<1$ the transform $\hat g(z)=\int e^{izx}g(x)\,dx$ is
\[
\hat g(z)=\int_{-\infty}^k e^{(1+iz)x}dx+e^k\int_k^\infty e^{izx}dx=\frac{e^{(1+iz)k}}{1+iz}-\frac{e^{(1+iz)k}}{iz}=\frac{e^{(1+iz)k}}{z^2-iz},
\]
and Parseval's identity gives $\E[g(x_T)]=\frac1{2\pi}\int\varphi(-z)\hat g(z)\,du$ along $z=u+\frac i2$. There $z^2-iz=u^2+\frac14$ and
$e^{(1+iz)k}=e^{k/2}e^{iuk}$. Changing $u$ to $-u$ and folding the integral onto $u>0$ gives the formula, since $Fe^{k/2}=\sqrt{FK}$.
\end{proposition}

The same argument prices any payoff whose transform is known on a strip, which is Lewis's point. A second route damps the call itself.

\begin{definition}[Carr--Madan formula]\label{def:dv:fourier-pricing-and-calibration-engineering:cm}
The \emph{Carr--Madan formula}\index{Carr--Madan formula} writes the damped call $e^{\alpha k}C(k)$, $\alpha>0$, as a Fourier integral in the
log-strike: with $F=1$,
\[
C(k)=\frac{e^{-\alpha k}}{\pi}\int_0^\infty\operatorname{Re}\bigl(e^{-ivk}\psi(v)\bigr)dv,\qquad
\psi(v)=\frac{P(0,T)\,\varphi\bigl(v-(\alpha+1)i\bigr)}{\alpha^2+\alpha-v^2+i(2\alpha+1)v},
\]
which on a grid in $v$ and $k$ is a discrete Fourier transform: the fast Fourier transform returns the calls on a whole grid of log-strikes
at once.
\end{definition}

The COS method (One Quant Book 4, chapter 28) takes a third route. It expands the density of $\ln(S_T/K)$ in a cosine series on a
truncated interval $[a,b]$. The coefficients come from $\varphi$ directly, and the put's payoff integrates against each cosine in closed form.
For smooth densities the error falls exponentially in the number of terms.

\begin{example}[Three pricers, one smile]\label{ex:dv:fourier-pricing-and-calibration-engineering:cos}
One-year calls in chapter 10's base \omterm{def:dv:stochastic-volatility:heston}{Heston model}, at 17 strikes from 60 to 140, against a reference (the Lewis integral by Gauss--Legendre
panels to $u=1\,000$, accurate to about $10^{-9}$; at the money it gives 6.7509). The COS method's largest error is 0.012 with 64 terms,
$8.6\times10^{-5}$ with 128 and $4.1\times10^{-7}$ with 256. Carr--Madan with interpolation between the grid's log-strikes reaches
$2.1\times10^{-4}$ with 1\,024 points, $1.8\times10^{-6}$ with 4\,096, and stalls near $2\times10^{-7}$
(\cref{fig:dv:fourier-pricing-and-calibration-engineering:transform}). The truncation interval matters as much as the terms. With $[a,b]$ at
twelve standard deviations instead of sixteen, COS stops at $3.0\times10^{-5}$ however many terms it gets, because Heston's left tail is
fatter than a normal's.
\end{example}

\begin{omfigure}
\begin{tikzpicture}
\begin{axis}[width=10.5cm,height=5cm,xmode=log,ymode=log,log basis x=2,
  xlabel={terms (COS) or grid points (Carr--Madan)},ylabel={largest price error},
  tick label style={font=\small},label style={font=\small},
  xmin=12,xmax=22000,ymin=1e-7,ymax=3,grid=major,grid style={gray!25},
  xtick={16,64,256,1024,4096,16384},xticklabels={16,64,256,1\,024,4\,096,16\,384},
  legend columns=2,legend cell align=left,legend style={font=\footnotesize,at={(0.5,-0.3)},anchor=north,draw=none,fill=none,/tikz/every even column/.append style={column sep=8pt}},
  table/col sep=comma]
\addplot[omThm,very thick,mark=*,mark size=1.6pt] table[x=n,y=err]{figdata/derivatives/24-fourier-pricing-and-calibration-engineering/cos.csv};
\addplot[omDef,very thick,mark=square*,mark size=1.6pt] table[x=n,y=err]{figdata/derivatives/24-fourier-pricing-and-calibration-engineering/fft.csv};
\legend{COS method,Carr--Madan FFT}
\end{axis}
\end{tikzpicture}

\omcaption{Largest error over 17 strikes of a one-year Heston call against the number of terms or grid points: COS converges exponentially to
the reference; the FFT converges fast and then stalls on its interpolation and damping errors. Data: the
tutorial.}\label{fig:dv:fourier-pricing-and-calibration-engineering:transform}
\end{omfigure}

For calibration the choice is practical. COS prices the strikes a desk quotes, where they are, in a few hundred operations each. The FFT
prices a fixed grid, and quoted strikes need interpolation. A smile of five strikes costs 800 characteristic-function evaluations by COS with 160 terms, and 4\,096 by the FFT.
A Levenberg--Marquardt calibration needs six residual evaluations per iteration (five for the Jacobian, one for the step) and, started from
yesterday, about eight iterations: some two hundred smiles.

\section{Calibration as a pipeline}

Book 4, chapter 24 treated calibration as an inverse problem. On a desk it is a job that runs every morning on data nobody has checked. The
build is a pipeline with five stages, each with its own failure modes.
\begin{enumerate}
\item \textbf{Quotes.} Bid and ask volatilities by expiry and strike, with the forward and the discount curve. A quote set is a snapshot,
kept, so that any past calibration can be rerun.
\item \textbf{Filter.} Drop quotes with no bid, crossed quotes (bid above ask), stale ones (spreads too wide to carry information) and
strikes too far out to matter. The synthetic market of this chapter corrupts one quote a day each way. The filter drops those two and keeps
the other eighteen.
\item \textbf{Objective.} Weighted errors, and optionally a penalty (next sections).
\item \textbf{Optimiser.} Levenberg--Marquardt, started from yesterday's parameters, in unconstrained coordinates
$z=(\ln v_0,\ln\kappa,\ln\bar v,\ln\eta,\operatorname{artanh}\rho)$ so that every step is admissible.
\item \textbf{Diagnostics.} Fit error overall and by expiry, the condition number of the Jacobian, parameter changes against limits. An alarm
stops the parameters from reaching the pricers until someone looks.
\end{enumerate}

\section{Objective, weights and constraints}

Calibrating in implied volatility is natural, since that is how errors are judged. But each residual then needs an implied-volatility
inversion. Calibrating in price is faster and needs a weight per quote.

\begin{definition}[Vega weighting]\label{def:dv:fourier-pricing-and-calibration-engineering:vega}
\emph{Vega weighting}\index{vega weighting} divides each price error by the Black \omterm{def:dv:greeks-and-the-hedging-pnl:greeks}{vega} of its quote, so that $(C^{\text{model}}-C^{\text{mid}})/
\mathcal V\approx\sigma^{\text{model}}-\sigma^{\text{mid}}$: a price calibration then minimises implied-volatility errors to first order,
without inverting.
\end{definition}

Without weights, price errors on long-dated at-the-money options dominate, since their \omterm{def:dv:greeks-and-the-hedging-pnl:greeks}{vegas} are largest.

\begin{example}[What the weights buy]\label{ex:dv:fourier-pricing-and-calibration-engineering:weights}
On the first day of the synthetic market (four expiries from one month to one year, five strikes each, eighteen after the filter), the Heston
fit's root-mean-square error by expiry, in volatility points, is:

\begin{center}\small
\begin{tabular}{lrrrrr}
\toprule
weights & 1 month & 3 months & 6 months & 1 year & all\\
\midrule
\omterm{def:dv:greeks-and-the-hedging-pnl:greeks}{vega} & 0.35 & 0.21 & 0.14 & 0.11 & 0.22\\
equal (price errors) & 0.46 & 0.15 & 0.12 & 0.06 & 0.24\\
\bottomrule
\end{tabular}
\end{center}

Equal weights buy a better one-year fit with a worse one-month one. At the money a one-month option's \omterm{def:dv:greeks-and-the-hedging-pnl:greeks}{vega} is the one-year's divided by
$\sqrt{12}$, so in price its errors count twelve times less.
\end{example}

Constraints come from the model and from the desk. The transformation keeps $v_0,\kappa,\bar v,\eta>0$ and $|\rho|<1$. The Feller condition
is not imposed. It fails in most equity calibrations, including this one ($2\kappa\bar v=0.19$ against $\eta^2=0.24$), and the simulation
schemes of chapter 23 are built for that case.

\section{Stability from day to day}

A calibration can reproduce the quotes perfectly and still be useless, if a small change in the quotes moves the parameters a lot. Cont and Ben
Hamida showed that the in-sample error can have many near-global minima, so that the calibration problem is ill-posed. The fit error then says
nothing about which minimum was picked.

\begin{example}[A flat valley]\label{ex:dv:fourier-pricing-and-calibration-engineering:valley}
The first day's fit is $v_0=0.0451$, $\kappa=1.82$, $\bar v=0.0518$, $\eta=0.488$, $\rho=-0.632$, with an error of 0.22 volatility point. Holding
$\eta$ fixed and refitting the other four gives errors of 0.29 at $\eta=0.35$, 0.23 at 0.45 and 0.55, and 0.27 at 0.65
(\cref{fig:dv:fourier-pricing-and-calibration-engineering:valley}, left). Across the whole range the error moves by less than the quotes'
noise of 0.3 volatility point, while $\kappa$ moves along with $\eta$. The market data (a \omterm{def:dv:jumps-and-levy-models:bates}{Bates model}, chapter 13, with noise) cannot tell these
\omterm{def:dv:stochastic-volatility:heston}{Heston models} apart.
\end{example}

\begin{omfigure}
\begin{tikzpicture}
\begin{groupplot}[group style={group size=2 by 1,horizontal sep=1.6cm},
  width=6.6cm,height=5cm,
  tick label style={font=\small},label style={font=\small},grid=major,grid style={gray!25},table/col sep=comma]
\nextgroupplot[xlabel={volatility of volatility $\eta$ (held fixed)},ylabel={fit error (vol points)},xmin=0.2,xmax=1.05,ymin=0,ymax=0.45,
  title={\small profile, first day}]
\addplot[omThm,very thick,mark=*,mark size=1.6pt] table[x=eta,y=rmse]{figdata/derivatives/24-fourier-pricing-and-calibration-engineering/profile.csv};
\addplot[gray,dashed,domain=0.2:1.05] {0.3};
\nextgroupplot[xlabel={fitted $\eta$},ylabel={fitted $\kappa$},xmin=0.2,xmax=0.8,ymin=0.5,ymax=4.5,title={\small 60 daily fits},
  legend columns=1,legend cell align=left,legend style={font=\footnotesize,at={(0.5,-0.32)},anchor=north,draw=none,fill=none}]
\addplot[omProp,only marks,mark=*,mark size=1.3pt] table[x=eta_free,y=kappa_free]{figdata/derivatives/24-fourier-pricing-and-calibration-engineering/daily.csv};
\addplot[omThm,only marks,mark=square*,mark size=1.3pt] table[x=eta_reg,y=kappa_reg]{figdata/derivatives/24-fourier-pricing-and-calibration-engineering/daily.csv};
\addplot[black,thick,dashed] table[x=eta,y=kappa]{figdata/derivatives/24-fourier-pricing-and-calibration-engineering/profile.csv};
\legend{no penalty,penalty $10^{-3}$,valley (profile)}
\end{groupplot}
\end{tikzpicture}

\omcaption{Left: the best fit error with the \omterm{def:dv:stochastic-volatility:sv}{volatility of volatility} held fixed, first day; the dashed line is the quotes' noise. Right: the
fitted $(\eta,\kappa)$ over 60 days without and with a penalty on parameter changes, and the first day's valley (the best $\kappa$ for each fixed
$\eta$): unpenalised, the fits scatter widely; penalised, they stay in a tight cluster.
Data: the tutorial.}\label{fig:dv:fourier-pricing-and-calibration-engineering:valley}
\end{omfigure}

Wandering along the valley costs money, because exotics are not priced by the fitted vanillas alone.

\begin{definition}[Recalibration P\&L]\label{def:dv:fourier-pricing-and-calibration-engineering:recalib}
The \emph{recalibration P\&L}\index{recalibration P\&L} of a position is the change in its model value caused by refitting the model's
parameters, with the market inputs that its hedges cover held fixed; in a P\&L attribution it lands in the unexplained column, and it has no
hedge.
\end{definition}

The remedy is to tell the optimiser that yesterday's parameters are informative. With $z$ the unconstrained parameters, the objective adds
$\lambda\lVert z-z_{\text{yesterday}}\rVert^2$ to the weighted squared errors, a Tikhonov penalty (Book 4, chapter 24). Along the valley the
penalty dominates and the parameters stay put. Across it the quotes dominate, and the parameters move when the market does.

\begin{example}[Sixty days with and without a penalty]\label{ex:dv:fourier-pricing-and-calibration-engineering:daily}
Unpenalised, the daily fits move $\eta$ by as much as 0.235 in a day, and the convexity adjustment of a one-year \omterm{def:dv:variance-swaps-and-volatility-derivatives:volswap}{volatility swap} (chapter 14)
changes by 0.36 volatility point a day (standard deviation). With $\lambda=10^{-3}$ the largest daily move of $\eta$ is 0.071, the
adjustment's daily changes fall to 0.07 volatility point, and the average fit error rises from 0.291 to 0.303 volatility point
(\cref{fig:dv:fourier-pricing-and-calibration-engineering:daily}). On day 59, where the free fit took $\eta$ from 0.51 to 0.28, the
penalised adjustment moves by 0.18 volatility point instead of 0.93.
\end{example}

\begin{omfigure}
\begin{tikzpicture}
\begin{groupplot}[group style={group size=2 by 1,horizontal sep=1.6cm},
  width=6.6cm,height=5cm,xlabel={day},xmin=0,xmax=61,
  tick label style={font=\small},label style={font=\small},grid=major,grid style={gray!25},table/col sep=comma]
\nextgroupplot[ylabel={volatility of volatility $\eta$},ymin=0.2,ymax=0.8,
  legend columns=2,legend cell align=left,legend style={font=\footnotesize,at={(1.12,-0.3)},anchor=north,draw=none,fill=none,/tikz/every even column/.append style={column sep=8pt}}]
\addplot[omProp,thick] table[x=day,y=eta_free]{figdata/derivatives/24-fourier-pricing-and-calibration-engineering/daily.csv};
\addplot[omThm,very thick] table[x=day,y=eta_reg]{figdata/derivatives/24-fourier-pricing-and-calibration-engineering/daily.csv};
\legend{no penalty,penalty $\lambda=10^{-3}$}
\nextgroupplot[ylabel={convexity adjustment (vol points)},ymin=0,ymax=2]
\addplot[omProp,thick] table[x=day,y=conv_free]{figdata/derivatives/24-fourier-pricing-and-calibration-engineering/daily.csv};
\addplot[omThm,very thick] table[x=day,y=conv_reg]{figdata/derivatives/24-fourier-pricing-and-calibration-engineering/daily.csv};
\end{groupplot}
\end{tikzpicture}

\omcaption{Sixty daily Heston calibrations to a noisy synthetic market. Left: the fitted \omterm{def:dv:stochastic-volatility:sv}{volatility of volatility}. Right: the convexity
adjustment of a one-year \omterm{def:dv:variance-swaps-and-volatility-derivatives:volswap}{volatility swap} (variance-swap volatility minus volatility-swap price) implied by each day's
fit. Data: the tutorial.}\label{fig:dv:fourier-pricing-and-calibration-engineering:daily}
\end{omfigure}

The penalty weight is a choice. Too small, and the parameters wander; too large, and they lag behind a market that moves. The build chooses
it by replaying history: for each $\lambda$ on a grid, run the daily calibration over past quote sets, record the largest daily parameter move
and the average fit error, and pick the smallest weight that meets the desk's limits
(\cref{fig:dv:fourier-pricing-and-calibration-engineering:scan}).

\begin{omfigure}
\begin{tikzpicture}
\begin{groupplot}[group style={group size=2 by 1,horizontal sep=2.2cm},
  width=6.4cm,height=5cm,xmode=log,xlabel={penalty weight $\lambda$},
  yticklabel style={/pgf/number format/fixed,/pgf/number format/precision=2},xmin=5e-6,xmax=2e-2,
  tick label style={font=\small},label style={font=\small},grid=major,grid style={gray!25},table/col sep=comma]
\nextgroupplot[ylabel={largest daily move of $\eta$},ymin=0,ymax=0.25]
\addplot[omThm,very thick,mark=*,mark size=1.6pt] table[x=reg,y=max_d_eta]{figdata/derivatives/24-fourier-pricing-and-calibration-engineering/scan.csv};
\addplot[gray,dashed,domain=5e-6:2e-2] {0.1};
\nextgroupplot[ylabel={rise in fit error (vol points)},ymin=-0.01,ymax=0.12]
\addplot[omDef,very thick,mark=square*,mark size=1.6pt] table[x=reg,y=d_rmse]{figdata/derivatives/24-fourier-pricing-and-calibration-engineering/scan.csv};
\addplot[gray,dashed,domain=5e-6:2e-2] {0.1};
\end{groupplot}
\end{tikzpicture}

\omcaption{Replaying 60 days for each penalty weight. Left: the largest daily move of the \omterm{def:dv:stochastic-volatility:sv}{volatility of volatility}, against a limit of 0.1.
Right: the rise in the average fit error, against a limit of 0.1 volatility point. Data: the
tutorial.}\label{fig:dv:fourier-pricing-and-calibration-engineering:scan}
\end{omfigure}

\section{Testing a calibration}

A calibration is code and gets the tests that code gets, plus three of its own.
\begin{itemize}
\item \textbf{Round trip.} Generate quotes from known parameters and recover them. From quotes made by $v_0=0.045$, $\kappa=2$, $\bar v=0.05$,
$\eta=0.5$, $\rho=-0.65$, the pipeline recovers every parameter to $10^{-11}$, from a start at chapter 10's base parameters. This tests the
pricer, the transformation and the optimiser together.
\item \textbf{Identification.} Add noise and refit many times. With 0.3 volatility point of noise, twenty refits give standard deviations of
0.0007 for $v_0$ (0.16 volatility point of at-the-money volatility), 0.049 for $\eta$, 0.046 for $\rho$, 0.0024 for $\bar v$ and 0.56 for
$\kappa$, 28\% of its value. The quotes pin down the short-term level and the skew, not the speed of mean reversion. The Jacobian's condition
number, 30 here, says the same in one number.
\item \textbf{Replay.} Rerun past quote sets with the production settings and check the parameter paths against the alarm limits, as in
\cref{fig:dv:fourier-pricing-and-calibration-engineering:scan}.
\end{itemize}
Only the round trip has a right answer. The other two measure what the quotes can support, which decides how exotics that depend on $\kappa$
should be reserved (chapter 27).

\section{Tutorial: transforms and a daily calibration}\label{tut:dv:fourier-pricing-and-calibration-engineering:tut}

\begin{tutorial}
\textbf{Goal.} Price a Heston smile by COS, Carr--Madan and the Lewis reference; calibrate a \omterm{def:dv:stochastic-volatility:heston}{Heston model} every day for 60 days of a noisy
synthetic market with and without a penalty on parameter changes, and choose the penalty weight by replay. \textbf{End state:} the four figures,
the weights table and the numbers of the weekend problem.
\begin{tutsteps}
\item \textbf{The COS pricer}, one row of cosine terms per strike:
\omcode{code/firm/calib/firm_calib.py}{30}{48}{Calls by the COS method.}\label{lst:dv:fourier-pricing-and-calibration-engineering:cos}
\item \textbf{The calibration}: vega-weighted residuals, the penalty on the change from yesterday, Levenberg--Marquardt, diagnostics:
\omcode{code/firm/calib/firm_calib.py}{164}{199}{Levenberg--Marquardt with a penalty on parameter changes.}\label{lst:dv:fourier-pricing-and-calibration-engineering:calib}
\item \textbf{Run} \texttt{dv\_calib.transform\_study()}, \texttt{weights\_table()}, \texttt{profile()}, \texttt{daily(0.0)}, \texttt{scan()},
\texttt{round\_trip()} and \texttt{fig\_calib.py}.
\end{tutsteps}
\textbf{What to change next.} Calibrate a \omterm{def:dv:jumps-and-levy-models:bates}{Bates model} to the same quotes and watch the jump parameters wander; penalise only $\kappa$ and $\eta$;
replace the Tikhonov penalty by a prior centred on a long-run average rather than on yesterday.
\end{tutorial}

\section{Build: the calibration pipeline}\label{bld:dv:fourier-pricing-and-calibration-engineering:calib}

\begin{build}
\bfield{Purpose} The miniature firm's calibration service: transform pricers for any model with a characteristic function, and a daily
pipeline that turns quotes into parameters the pricers can trust.
\bfield{Interface} \texttt{cos\_calls(model, fwd, strikes, t, n, width)}, \texttt{carr\_madan(model, fwd, t, n, eta, alpha)},
\texttt{carr\_madan\_at}, \texttt{lewis\_calls} (reference), \texttt{cumulants}; \texttt{QuoteSet(t, strike, bid, ask, fwd)},
\texttt{filter\_quotes(q, max\_spread, max\_sd)}; \texttt{calibrate(q, z0, make, weights, prior, reg)} $\to$ \texttt{Fit(model, z, rmse,
penalty, iterations, condition)}; \texttt{vol\_errors}, \texttt{jump\_alarm(z\_new, z\_old, limits)}; \texttt{heston\_from}, \texttt{heston\_to}.
\bfield{Rules} Quote sets are stored and every calibration can be replayed; the filter reports what it dropped and why; weights are explicit;
the calibration starts from yesterday's parameters; an alarm blocks parameters that move beyond their limits.
\bfield{Acceptance tests} \texttt{code/firm/calib/tests/}: COS, Carr--Madan and the Lewis reference against Black's formula and against each other;
the cumulants of a lognormal; the filter's four reasons; the round trip; an infinite penalty keeps yesterday's parameters; the vega-weighted
price error equals the volatility error to first order.
\bfield{Stretch} The FFT engine for the time-dependent models of chapter 12; analytic Jacobians from the characteristic function's derivatives;
a calibration of the \omterm{def:dv:jumps-and-levy-models:bates}{Bates model} with the jump parameters penalised to a long-run prior.
\end{build}

\begin{omsources}
\item P.~Carr and D.~B.~Madan, ``Option valuation using the fast Fourier transform'', \emph{Journal of Computational Finance} 2(4) (1999) 61--73.
\item A.~L.~Lewis, ``A simple option formula for general jump-diffusion and other exponential L\'evy processes'', working paper (2001).
\item F.~Fang and C.~W.~Oosterlee, ``A novel pricing method for European options based on Fourier-cosine series expansions'', \emph{SIAM Journal on
Scientific Computing} 31(2) (2008) 826--848.
\item R.~Cont and S.~Ben Hamida, ``Recovering volatility from option prices by evolutionary optimization'', \emph{Journal of Computational Finance}
8(4) (2005) 43--76.
\end{omsources}

\section{Exercises}

\begin{exercise}[$\star$]\label{exo:dv:fourier-pricing-and-calibration-engineering:1}
Check that the \omterm{def:dv:fourier-pricing-and-calibration-engineering:lewis}{Lewis formula} gives $C=P(0,T)F$ at $K=0$ and that $\varphi(-i)=1$ is what makes $x_T$ a forward-martingale log-price.
\end{exercise}

\begin{exercise}[$\star$]\label{exo:dv:fourier-pricing-and-calibration-engineering:2}
Why does the \omterm{def:dv:fourier-pricing-and-calibration-engineering:cm}{Carr--Madan formula} need the damping factor $e^{\alpha k}$, and what limits $\alpha$ from above?
\end{exercise}

\begin{exercise}[$\star$]\label{exo:dv:fourier-pricing-and-calibration-engineering:3}
Show that a price error divided by the Black \omterm{def:dv:greeks-and-the-hedging-pnl:greeks}{vega} is the implied-volatility error to first order.
\end{exercise}

\begin{exercise}[$\star\star$]\label{exo:dv:fourier-pricing-and-calibration-engineering:4}
With 0.3 volatility point of quote noise and twenty quotes, what fit error should a correct model reach? Compare with the chapter's fits.
\end{exercise}

\begin{exercise}[$\star\star$]\label{exo:dv:fourier-pricing-and-calibration-engineering:5}
Why does COS stall at $3\times10^{-5}$ with a twelve-standard-deviation interval, and why do more terms not help?
\end{exercise}

\begin{exercise}[$\star\star$]\label{exo:dv:fourier-pricing-and-calibration-engineering:6}
Explain why $\kappa$ and $\eta$ move together along the valley of \cref{fig:dv:fourier-pricing-and-calibration-engineering:valley}.
\end{exercise}

\begin{exercise}[$\star\star\star$]\label{exo:dv:fourier-pricing-and-calibration-engineering:7}
\emph{Coding.} Rerun the 60-day calibration with equal price weights instead of \omterm{def:dv:greeks-and-the-hedging-pnl:greeks}{vega} weights, unpenalised. Does the \omterm{def:dv:stochastic-volatility:sv}{volatility of volatility}
wander more or less, and why?
\end{exercise}

\begin{exercise}[$\star\star\star$]\label{exo:dv:fourier-pricing-and-calibration-engineering:8}
\emph{Find the flaw.} ``Our calibration is excellent: the fit error has been below 0.45 volatility point every day for two months.''
\end{exercise}

\section{Problem: The Wandering Parameter}

\begin{problem}[{Weekend problem --- how much to penalise}]\label{pb:dv:fourier-pricing-and-calibration-engineering:1}
A desk calibrates a \omterm{def:dv:stochastic-volatility:heston}{Heston model} every morning to twenty quotes (one, three, six and twelve months; five strikes each) and prices \omterm{def:dv:variance-swaps-and-volatility-derivatives:volswap}{volatility swaps}
and \omterm{def:dv:asians-lookbacks-cliquets-and-forward-starts:cliquet}{cliquets} with it. Over sixty days the unpenalised \omterm{def:dv:stochastic-volatility:sv}{volatility of volatility} wanders.

\medskip\noindent\textbf{Part I --- The pricer.}
\begin{enumerate}
\item Which of COS and Carr--Madan would you use for this calibration, and why?
\item How many COS terms give a price error below $10^{-4}$ on the one-year smile?
\item What sets the truncation interval, and what happens when it is too narrow?
\item Roughly how long does one calibration take, and what dominates it?
\item How would you check the pricer independently of the calibration?
\end{enumerate}

\medskip\noindent\textbf{Part II --- The fit.}
\begin{enumerate}[resume]
\item What does the filter drop on a typical day in the synthetic market?
\item Give the first day's parameters and fit error.
\item What happens to the one-month fit with equal price weights?
\item How flat is the fit error along $\eta$, compared with the noise?
\item What does the round trip with noise say about which parameters the quotes identify?
\end{enumerate}

\medskip\noindent\textbf{Part III --- Day to day.}
\begin{enumerate}[resume]
\item Give the largest daily move of $\eta$ without a penalty, and the day it happens.
\item What does it do to the \omterm{def:dv:variance-swaps-and-volatility-derivatives:volswap}{volatility swap}'s convexity adjustment, and what is that in money for a \omterm{def:dv:variance-swaps-and-volatility-derivatives:veganot}{vega notional} of 100\,000 per volatility point?
\item Why is this move a \omterm{def:dv:fourier-pricing-and-calibration-engineering:recalib}{recalibration P\&L} and not a market move?
\item How does the penalty change the parameter and the adjustment paths?
\item What does the penalty cost in fit error?
\end{enumerate}

\medskip\noindent\textbf{Part IV --- Choosing the weight.}
\begin{enumerate}[resume]
\item Describe the replay that chooses $\lambda$.
\item What is the risk of a penalty that is too large?
\item Which parameters would you penalise more, and which less?
\item State the \emph{named result}: the regularisation weight at which the day-to-day moves of the \omterm{def:dv:stochastic-volatility:sv}{volatility of volatility} fall below 0.1 while the
fit error rises by less than 0.1 volatility point.
\item In one sentence: what does a penalty on parameter changes buy?
\end{enumerate}
\end{problem}

\section{Interview questions}

\begin{interviewq}[$\star$ \iqroles{developer, researcher}]\label{iq:dv:fourier-pricing-and-calibration-engineering:1}
How do you price a European option when you only know the characteristic function?
\end{interviewq}

\begin{interviewq}[$\star\star$ \iqroles{researcher}]\label{iq:dv:fourier-pricing-and-calibration-engineering:2}
Compare the Carr--Madan FFT and the COS method.
\end{interviewq}

\begin{interviewq}[$\star\star$ \iqroles{researcher, bank}]\label{iq:dv:fourier-pricing-and-calibration-engineering:3}
Your Heston calibration's parameters jump from day to day while the fit stays good. What is happening and what do you do?
\end{interviewq}

\begin{interviewq}[$\star\star$ \iqroles{developer}]\label{iq:dv:fourier-pricing-and-calibration-engineering:4}
Design a daily calibration job. What does it store, check and alarm on?
\end{interviewq}

\begin{interviewq}[$\star\star$ \iqroles{trader, risk}]\label{iq:dv:fourier-pricing-and-calibration-engineering:5}
Your book shows a large unexplained P\&L on a quiet day. How could the calibration be the cause?
\end{interviewq}

\begin{interviewq}[$\star\star\star$ \iqroles{researcher}]\label{iq:dv:fourier-pricing-and-calibration-engineering:6}
Why should you weight a calibration, and how would you choose the weights?
\end{interviewq}
